% ===== §1 =====
\section{The Limit Proposed, Not Lamented}
\label{p26:sec:intro}

\noindent
There is an old habit of treating the limit of language as a defect. On this habit, language is an instrument for the transport of meaning, and wherever it fails to carry meaning across, it has fallen short of what it was for. The failure is then met with resignation, with the wish for a fuller language, or with a mysticism that points past speech to a silence it cannot itself enter. This paper takes a different route. It does not lament the limit of language and does not propose to overcome it. It proposes that the limit is the condition of something, and that the something is the very life of relation.

\subsection*{The proposal}

The proposal can be stated at once. In each of the dimensions along which a relation lives, its coming into being, its being kept alive, the value it generates, the ethics it asks of those within it, the justice it owes, and the way it grows and turns into new forms, language reaches toward something it does not reach. This paper first lays out these six limits one by one. It gives each of them its own shape, so that the limit of language is seen not as a single vague insufficiency but as six distinct and structured phenomena, each with its own core that speech circles and does not touch, and each stated as a claim of its own.

The paper then makes its central move. What stands at the limit, in each of the six cases, is not an absence to be regretted. It is the site at which drive, in the sense psychoanalysis gave the word, becomes possible. Drive is not one more content that language happened to leave out and might, with better instruments, recover. Drive is the phenomenology of that limit itself. It is the way the gap between what can be said and what presses beneath saying shows itself, the felt form of a remainder that the operation of the signifier can never close. This is why the attempt to say the drive returns not the drive but its renewed escape. To symbolize the gap is to find that the gap has moved, because the gap is nothing other than what symbolization cannot close.

\subsection*{What follows from it}

From this the rest follows. If the limit is the site of drive, then the non-closure of language is not a wound in relation but the open space in which relation keeps moving. A relation in which everything had been said, in which the other had been wholly reached, in which no remainder stirred beneath the words, would be a relation in which nothing further pressed, nothing further turned, nothing further was generated. The completeness so often wished for would be the end of the very thing it wished to secure. This consequence, argued in Part Two by a thought experiment, governs the whole account: closure is not fulfillment but extinction, and the openness of the limit is what keeps a bond alive.

\subsection*{The path of the paper}

The paper proceeds in five parts. Part One sets out the six limits, giving each its own structure and its own claim: the ungroundability of the bond's existence, the uncapturability of the living present, the absence of any external place to price relational value, the aesthetic measure of the ethical response, the generativity that justice guards, and the direction of growth that cannot be said in advance. Part Two traces the shared form of these six to its source in the constitution of the speaking subject, gives the gap its two faces and its object, distinguishes the drive that circles from the desire that slides, and establishes by a thought experiment that the closure of the gap would extinguish rather than fulfill them, from which the paper's central claim follows, that drive is the phenomenology of the limit. Part Three follows the persistence of drive into the time of relation, where an early signifier goes on organizing later bonds, the symptom appears as a fidelity to what cannot be said, and the good bond is shown to consist in making room for the other's unsymbolizable core. Part Four draws out the generativity of the limit in three registers, the poem as a language that circles the gap without closing it, an ethics that generates a bond without possessing the other, and a justice that guards the generativity of relational being against premature settlement. Part Five closes without a synthesis, in keeping with what the paper argues, since a work about what language cannot reach ought not to pretend, in its last pages, to have reached it.

\subsection*{Method and stance}

Two remarks on method. First, the paper states its own contributions as named claims, set off from the prose, so that what it proposes can be found and weighed apart from the argument that leads to it. Second, at each dimension and each register the paper passes through the relevant literature not to survey it but to locate, in each account, the point it reaches toward and the remainder it leaves, so that the review is everywhere in the service of the position and never a catalogue.

A remark on stance. This is a constructive paper and not a polemical one. It sets up no adversary and refutes no school. Where it passes near positions that would treat full communication as an ideal, or transparency between persons as the measure of a good bond, it does not stop to defeat them but builds a different account and lets the difference speak. The traditions it engages, in speech-act theory, in phenomenology, in the theory of value, in psychoanalysis, in the ethics of the other, and in the study of testimony, are taken as sources of insight to be positioned, each having reached toward the same limit from its own side.
